\documentclass[sigconf, nonacm]{acmart}

\acmConference[ACM-BCB 2026]{The 17th ACM Conference on Bioinformatics,
  Computational Biology, and Health Informatics}{June 30, 2026}{Calabria, Italy}
\setcopyright{none}
\acmDOI{}
\acmISBN{}
\settopmatter{printacmref=false, printccs=false, printfolios=true}

% Aggressive spacing trim to fit 2 pages
\setlength{\textfloatsep}{4pt}
\setlength{\floatsep}{4pt}
\setlength{\intextsep}{4pt}
\setlength{\abovecaptionskip}{2pt}
\setlength{\belowcaptionskip}{1pt}

\usepackage{booktabs}
\usepackage{amsmath}
\usepackage{microtype}
\usepackage{url}

\begin{document}

\title[Protocol Choices Dominate GRN Benchmarking]{Evaluation Protocol Choices Dominate Gene Regulatory Network Benchmarking Outcomes for Single-Cell Foundation Models}

\author{Ihor Kendiukhov}
\email{kenduhov.ig@gmail.com}
\affiliation{%
  \institution{Independent Researcher}
  \country{Ukraine}
}
\renewcommand{\shortauthors}{Kendiukhov}

\begin{abstract}
Recent benchmarks report that single-cell foundation models such as scGPT
and Geneformer recover gene regulatory networks (GRNs) from attention
weights at performance competitive with established baselines. We show
that this conclusion is governed primarily by three evaluation-protocol
choices---gene-symbol mapping, the candidate edge set, and the gold
standard---rather than by the models. We introduce the
\emph{universe-aware} protocol that scores every pair in
$U = T \times (G \setminus \{\text{self}\})$ where $T$ is a pinned
human-TF list and $G$ is the set of HGNC-approved expressed genes, and
apply it to BEELINE hESC and hHep with cell-type ChIP-seq,
non-specific ChIP-seq, and STRING gold standards. Across 279 (dataset,
gold standard, candidate set, method) tuples, an ANOVA decomposition of
$\log_{10}(\text{AUPR})$ assigns $80.1\%$ of the variance to
protocol-related factors and $0.27\%$ to method identity. Restricting
the candidate set to the literature-standard
$T \times T_{\text{gold-targets}}$ inflates AUPR by $13\text{--}17\times$
without changing the underlying scores; AUPR tracks base rate at
$\rho = 0.997$. Under universe-aware scoring against cell-type ChIP-seq,
all six scGPT attention-aggregation variants have AUROC below $0.5$ on
hHep (95\% DeLong CI strictly below $0.5$, $p \ll 0.001$), while the
same scGPT predictions reach AUROC $0.58$ on hESC; the verdict flips by
dataset and is invisible under aggregated single-number benchmarks. We
release a versioned pipeline and a six-item reproducibility checklist.
Code: \url{https://github.com/Biodyn-AI/grn-evaluation-bias}.
\end{abstract}

\keywords{gene regulatory networks, single-cell foundation models,
benchmarking, scGPT, Geneformer, reproducibility}

\maketitle

\section{Universe-Aware Evaluation}

For a single-cell dataset $D$ let $G(D)$ be the HGNC-approved expressed
genes after BEELINE quality control under a SHA-pinned symbol
dictionary, and let $T(D) = T_{\text{BEELINE}} \cap G(D)$ be the
intersection of the 1{,}563-entry BEELINE human-TF list with $G(D)$.
The \emph{universe-aware} candidate set is
\[
U(D) = \{ (s,t) : s \in T(D),\ t \in G(D),\ s \neq t \}.
\]
For BEELINE \textbf{hESC} (17{,}735 genes $\times$ 758 cells, Chu et al.~2016):
$|G|{=}17{,}724$, $|T|{=}1{,}459$, $|U|{=}25{,}857{,}857$;
\textbf{hHep} (11{,}515 $\times$ 425, Camp et al.~2017):
$|G|{=}11{,}512$, $|T|{=}1{,}095$, $|U|{=}12{,}604{,}545$. All methods
score \emph{every} pair in $U(D)$.

\textbf{Methods compared.} scGPT (whole-human, 12 layers, 8 heads)
and Geneformer-V1-10M (6 layers, 4 heads) attention with six
aggregation variants; baselines Pearson, Spearman, GRNBoost2 (100
trees), and a seeded random scorer. AUROC significance uses DeLong's
analytical SE~\cite{delong1988}; AUPR difference uses 2{,}000-resample
paired bootstrap.

\section{Results}

\textbf{Protocol dominates model.}
A Type-II ANOVA on $\log_{10}(\text{AUPR})$ across 279
(dataset, gold standard, candidate set, method) tuples assigns
$57.5\%$ of variance to gold standard and $22.6\%$ to candidate set
(cumulative $80.1\%$ protocol), versus $0.27\%$ to method identity---a
$\sim$300$\times$ gap. Restricting the candidate set from $T \times G$
to $T \times T_{\text{gold-targets}}$ inflates median AUPR by
$12.6\times$ on hESC and $17.0\times$ on hHep with the same scores;
inflation is uniform across methods including the random baseline.
Across all tuples, Pearson correlation between AUPR and base rate is
$\rho = 0.997$.

\textbf{Foundation-model verdict flips by dataset.}
Table~\ref{tab:headline} shows universe-aware AUPR and AUROC against
cell-type-specific ChIP-seq. On hHep, all six scGPT aggregation
variants have AUROC $\le 0.499$ with $95\%$ analytical CI
upper bound below $0.5$ ($\sigma_{\text{AUROC}} \le 6\!\times\!10^{-4}$,
$p \ll 0.001$ vs $H_0: \text{AUROC}=0.5$), i.e., significantly
\emph{below random}; on hESC the same model reaches AUROC $\sim 0.58$.
Geneformer is moderate on hHep (AUROC $0.53$) and strong on hESC
(AUROC $0.64$). On the broader STRING gold standard, no method is
below random. The dataset $\times$ gold-standard sensitivity is
invisible in aggregated benchmarks that report a single AUPR number,
and the aggregation flatters the model because it averages
``below-random on hHep'' with ``above-random on hESC''.

\begin{table}[t]
\caption{Universe-aware AUPR and AUROC, cell-type ChIP-seq.
Foundation models: median (range) across six aggregations.
$\dagger$: 95\% DeLong CI for AUROC lies entirely below $0.5$.}
\label{tab:headline}
\footnotesize
\setlength{\tabcolsep}{4pt}
\begin{tabular}{llcc}
\toprule
Dataset & Method & AUPR & AUROC \\
\midrule
hHep        & Spearman                & 0.020 & 0.526 \\
$b{=}0.018$ & Geneformer (md.)        & 0.020 & 0.531\,(.529-.534) \\
$|U|{=}12.6$M & Pearson               & 0.019 & 0.518 \\
            & \emph{random}           & 0.018 & 0.501 \\
            & GRNBoost2               & 0.018 & 0.506 \\
            & \textbf{scGPT (md.)}    & \textbf{0.018} & \textbf{0.497$^{\dagger}$}\,(.493-.499) \\
\midrule
hESC        & Geneformer (md.)        & 0.020 & 0.638\,(.635-.640) \\
$b{=}0.013$ & scGPT (md.)             & 0.017 & 0.578\,(.577-.580) \\
$|U|{=}25.9$M & Pearson               & 0.015 & 0.550 \\
            & Spearman                & 0.015 & 0.544 \\
            & GRNBoost2               & 0.014 & 0.522 \\
            & \emph{random}           & 0.013 & 0.500 \\
\bottomrule
\end{tabular}
\end{table}

\textbf{Aggregation choice does not rescue the result.} All six
aggregations fall within a tight range on universe-aware AUPR
(scGPT hHep: $0.0183$--$0.0187$; Geneformer hHep:
$0.0195$--$0.0199$); the ``max-heads, mean-layers'' variant cited in
prior reports is statistically indistinguishable from the
mean-by-mean default. The full Type-II decomposition is
Table~\ref{tab:anova}.

\begin{table}[t]
\caption{Type-II sum of squares of $\log_{10}(\text{AUPR})$ across
279 (dataset, gold standard, candidate set, method) tuples. Protocol
factors account for $80.1\%$; method identity for $0.27\%$.}
\label{tab:anova}
\footnotesize
\setlength{\tabcolsep}{6pt}
\begin{tabular}{lrr}
\toprule
Factor & SS & \% total \\
\midrule
Gold standard           & 42.88 & \textbf{57.5\%} \\
Candidate set           & 16.89 & \textbf{22.6\%} \\
Dataset (hESC vs hHep)  & 0.42  & 0.6\% \\
Method family           & 0.20  & \textbf{0.27\%} \\
Residual                & 14.13 & 18.9\% \\
\bottomrule
\end{tabular}
\end{table}

\textbf{Symbol mapping shifts coverage.}
Under the pinned HGNC dictionary, $89$--$93\%$ of raw symbols resolve
as approved, $7$--$10\%$ via previous/alias resolution, and $\sim 1\%$
are unmapped. Without HGNC mapping, edge coverage between
predicted edges and the cell-type ChIP-seq gold is effectively zero;
with mapping it exceeds $99\%$. Symbol-mapping choices are therefore
necessary upstream of every numeric claim in this and prior reports;
we ship the dictionary as a SHA-checked file.

\textbf{Gold-standard sensitivity is the third axis of protocol bias.}
The same scGPT predictions move from \emph{below random} on hHep
cell-type ChIP-seq to clearly above random on STRING and on hESC
cell-type ChIP-seq (Table~\ref{tab:gs_sens}). Geneformer similarly
shifts by $0.07$--$0.11$ AUROC across gold standards on the same
dataset. Reports that cite a single ``AUPR against ChIP-seq'' number
for a foundation model are not falsifiable without disclosing the
gold-standard identity, the candidate set, and the universe size.

\begin{table}[t]
\caption{Universe-aware AUROC across three gold standards, median
across attention-aggregation variants. The same model gives
\emph{opposite verdicts} (below vs above random, $\dagger = $ 95\%
DeLong CI below $0.5$) depending on the gold standard.}
\label{tab:gs_sens}
\footnotesize
\setlength{\tabcolsep}{4pt}
\begin{tabular}{llccc}
\toprule
& & cell-type & non-specific & \\
Tag & Method & ChIP-seq & ChIP-seq & STRING \\
\midrule
hHep & scGPT       & 0.497$^\dagger$ & 0.476$^\dagger$ & 0.579 \\
     & Geneformer  & 0.531           & 0.499           & 0.595 \\
     & GRNBoost2   & 0.506           & 0.501           & 0.543 \\
     & Pearson     & 0.518           & 0.508           & 0.555 \\
     & Spearman    & 0.526           & 0.519           & 0.590 \\
\midrule
hESC & scGPT       & 0.578           & 0.523           & 0.604 \\
     & Geneformer  & 0.638           & 0.550           & 0.665 \\
     & GRNBoost2   & 0.522           & 0.501           & 0.544 \\
     & Pearson     & 0.550           & 0.517           & 0.585 \\
     & Spearman    & 0.544           & 0.513           & 0.575 \\
\bottomrule
\end{tabular}
\end{table}

\textbf{Implications for benchmark design.}
A defensible foundation-model GRN claim must specify three artifacts
along with every metric: the pinned symbol dictionary, the candidate
universe $U$ (and its cardinality and base rate), and the gold standard
with its source/cell-type/version. The literature today routinely
specifies only the last of these three, and often only loosely. Our
released pipeline emits all three by construction for every
evaluation; we hope it lowers the cost of reporting these specifics
in future work.

\textbf{Top-of-ranking precision is even harsher than AUROC.}
At the top of the universe-aware ranking on hHep cell-type ChIP-seq
(base rate $0.018$, so random expects $\approx 0.9$ positives in the
top $50$), precision$@50$ is zero for $5$ of $6$ scGPT variants and
for Pearson and Spearman; precision$@1000$ is $0.009$--$0.017$ for
scGPT versus $0.026$ for random. Geneformer fares better
(precision$@50 = 0.04$--$0.06$, precision$@1000 = 0.025$--$0.028$,
$\sim$$1.0$--$1.6\times$ random). For both foundation models on hHep,
the top of the ranked attention edge list is statistically
indistinguishable from---or worse than---a random shuffle.

\section{Reproducibility Checklist}

(1) pin and ship a SHA-checked gene-symbol dictionary
(here: HGNC dump SHA-256 \texttt{b9482a52\ldots}, 50{,}094 rows);
(2) define and report $|G|$, $|T|$, $|U|$, and base rate for every
metric;
(3) report AUPR with its candidate-set definition---never as a single
number---and report AUROC alongside;
(4) declare ``below random'' only when an analytical or
bootstrap-derived $95\%$ CI for the relevant statistic excludes the
null;
(5) include both a correlation baseline (Pearson/Spearman) and a
tree-based baseline (GRNBoost2/GENIE3) under the same universe;
(6) release executable pipelines with frozen dependency versions,
seeded permutation/bootstrap, and SHA-checked data archives
(here: BEELINE Zenodo record 3701939, SHA \texttt{1fc84949\ldots} for
inputs, \texttt{d54fc6ed\ldots} for networks).

\textbf{Limitations.} We evaluate two human BEELINE datasets;
mouse datasets and other GRN benchmarks (e.g.,~DREAM5) are left to
future work. scGPT's architectural sequence cap of $1{,}200$ tokens
bounds the per-cell context, but cell-aggregated attention is
collected over all expressed genes; results may differ for the
$316$M-parameter Geneformer-V2 we did not run. BEELINE expression is already log-normalized; we invert via
$\mathrm{round}(\mathrm{expm1}(\cdot))$ for the tokenizer-based
models and ship the inversion script with the pipeline.

Our analysis does \emph{not} preclude the possibility that
foundation-model representations contain regulatory signal extractable
by more sophisticated decoders; it establishes that raw attention
aggregation does not suffice and that current benchmark claims require
re-examination under controlled evaluation protocols.

\textbf{Code and data availability.}
The full pipeline---BEELINE preprocessing, attention extraction for
scGPT and Geneformer, baselines, universe-aware evaluation,
DeLong AUROC CI with paired AUPR bootstrap, Type-II ANOVA
decomposition, and a paper-snippet generator that traces every
number in this abstract back to its source CSV---is released under
MIT at \url{https://github.com/Biodyn-AI/grn-evaluation-bias},
together with the pinned HGNC dictionary, the BEELINE SHA checks,
and a conda lockfile. All numeric values in this work are emitted
verbatim by the released pipeline and are independently re-derivable.

\textbf{Acknowledgments.} We thank the BEELINE, scGPT, and Geneformer
authors for releasing their data and checkpoints openly. No external
funding supported this work.

\begin{thebibliography}{1}
\footnotesize
\bibitem{cui2024scgpt}
H.~Cui et al.
sc{GPT}: toward building a foundation model for single-cell
multi-omics using generative {AI}.
\emph{Nature Methods}, 21(8):1470--1480, 2024.

\bibitem{theodoris2023geneformer}
C.~V. Theodoris et al.
Transfer learning enables predictions in network biology.
\emph{Nature}, 618(7965):616--624, 2023.

\bibitem{pratapa2020}
A.~Pratapa et al.
Benchmarking algorithms for gene regulatory network inference from
single-cell transcriptomic data.
\emph{Nature Methods}, 17(2):147--154, 2020.

\bibitem{moerman2019grnboost2}
T.~Moerman et al.
GRNBoost2 and Arboreto: efficient and scalable inference of gene
regulatory networks. \emph{Bioinformatics}, 35(12):2159--2161, 2019.

\bibitem{delong1988}
E.~R. DeLong, D.~M. DeLong, and D.~L. Clarke-Pearson.
Comparing the areas under two or more correlated ROC curves: a
nonparametric approach. \emph{Biometrics}, 44(3):837--845, 1988.

\bibitem{tweedie2021hgnc}
S.~Tweedie et al.
Genenames.org: the HGNC and VGNC resources in 2021.
\emph{Nucleic Acids Research}, 49(D1):D939--D946, 2021.

\bibitem{lambert2018tfs}
S.~A. Lambert et al.
The human transcription factors. \emph{Cell}, 172(4):650--665, 2018.

\bibitem{string2019}
D.~Szklarczyk et al.
STRING v11: protein--protein association networks with increased
coverage. \emph{Nucleic Acids Research}, 47(D1):D607--D613, 2019.
\end{thebibliography}

\end{document}
